% ============================================================================
%  MATEMÁTICAS OPERATIVAS — Sesión 15
%  Semana 15: 16 al 21 de noviembre de 2026
%  (lunes 16: festivo, Independencia de Cartagena)
%  3.5 Ecuaciones de primer grado · 3.6 Ecuaciones cuadráticas
%  Repaso integrador de la Unidad 3 y simulacro del Examen 3
%  Duración prevista: 3 horas
% ============================================================================
\documentclass[aspectratio=169,10pt]{beamer}

\usepackage[utf8]{inputenc}
\IfFileExists{spanish.ldf}{\usepackage[spanish,es-noshorthands]{babel}}{}
\usepackage{amsmath,amssymb}
\usepackage{booktabs}
\usepackage{tikz}

\usetheme{Uniminuto}

\uniasignatura{Matemáticas Operativas}
\uniprofesor{Andrés Felipe Puerta Velez}
\unicorreo{andres.puerta-v@uniminuto.edu.co}
\uninrc{CBASBL011}
\unisemana{Sesión 15 — Semana 15 · 16 al 21 de noviembre de 2026}

\newcommand{\listatiny}{\setbeamerfont{itemize/enumerate body}{size=\tiny}\tiny}

\begin{document}

\portada

\begin{frame}{Orden del día}
  \begin{enumerate}
    \item Revisión de los ejercicios propuestos de la sesión 14
    \item \azul{Tema 3.5: ecuaciones de primer grado y aplicaciones}
    \item \azul{Tema 3.6: ecuaciones cuadráticas y aplicaciones}
    \item Repaso integrador de la Unidad 3
    \item Simulacro del Examen 3 y recomendaciones
  \end{enumerate}

  \vspace{0.25cm}
  {\footnotesize\color{uniGris} Última sesión de clase del semestre.
  \destacar{Lunes 16 fue festivo} (Independencia de Cartagena). Hoy cerramos
  los dos temas que faltan: el \destacar{Examen 3 (30\%)} se presenta la
  próxima semana en Aulas Virtuales.}
\end{frame}

% ===========================================================================
\seccion{3.5 Ecuaciones de primer grado}
% ===========================================================================

\begin{frame}{Resolver es aislar la incógnita}
  \scriptsize
  \begin{block}{El principio}
    \tiny
    Una ecuación es una \azul{balanza}: lo que se le haga a un lado hay que
    hacérselo al otro. Se puede sumar, restar, multiplicar y dividir por lo
    mismo —\destacar{salvo por cero}— sin cambiar las soluciones.
  \end{block}

  \vspace{0.04cm}
  \begin{columns}[T,onlytextwidth]
    \begin{column}{0.48\textwidth}
      \begin{exampleblock}{Ejemplo básico}
        \tiny
        \[ 5(x-2) = 3x+8 \]
        \[ 5x-10 = 3x+8 \]
        \[ 5x-3x = 8+10 \;\Longrightarrow\; 2x = 18 \]
        \[ x = \destacar{9} \]
        \azul{Verif.:} $5(7)=35$ y $3(9)+8=35$ \checkmark
      \end{exampleblock}
    \end{column}
    \begin{column}{0.48\textwidth}
      \begin{exampleblock}{Con fracciones: quite denominadores}
        \tiny
        \[ \frac{x-1}{2}+\frac{x+3}{5} = 4 \]
        Multiplique todo por el m.c.d. $=10$:
        \[ 5(x-1)+2(x+3) = 40 \]
        \[ 7x+1 = 40 \;\Longrightarrow\; x = \destacar{\tfrac{39}{7}} \]
      \end{exampleblock}
    \end{column}
  \end{columns}

  \vspace{0.0cm}
  {\tiny\color{uniGris} \destacar{El orden que menos errores produce:}
  (1) quitar denominadores, (2) quitar paréntesis, (3) pasar las $x$ a un lado
  y los números al otro, (4) despejar, (5) \azul{verificar}. La verificación
  no es opcional: es gratis.}
\end{frame}

\begin{frame}{Traducir un problema a una ecuación}
  \scriptsize
  \begin{center}
  \begin{tabular}{@{}p{5.6cm}p{6.4cm}@{}}
    \toprule
    \textbf{En palabras} & \textbf{En símbolos} \\
    \midrule
    Un número aumentado en 7 & $x+7$ \\
    El doble de un número, menos 5 & $2x-5$ \\
    Tres números enteros consecutivos & $x$, \ $x+1$, \ $x+2$ \\
    El $19\%$ más sobre un precio & $1{,}19x$ \\
    La edad dentro de 4 años & $x+4$ \\
    Un rectángulo $7$ m más largo que ancho & ancho $x$, largo $x+7$ \\
    \bottomrule
  \end{tabular}
  \end{center}

  \vspace{0.02cm}
  \begin{exampleblock}{Ejemplo desarrollado}
    \tiny
    Un artículo con IVA del $19\%$ cuesta $\$297\,500$. ¿Cuál es el precio
    sin IVA? \ $1{,}19x = 297\,500 \Rightarrow
    x = \tfrac{297\,500}{1{,}19} = \destacar{\$250\,000}$.
    \ \azul{Verificación:} $250\,000\cdot 0{,}19 = 47\,500$ y
    $250\,000+47\,500 = 297\,500$ \checkmark
  \end{exampleblock}

  \vspace{0.02cm}
  {\tiny\color{uniGris} \destacar{Lo más difícil de un problema no es
  resolverlo, es plantearlo.} Empiece siempre escribiendo \azul{``sea $x$ =
  \dots''} con toda claridad; la mitad de los errores vienen de no saber qué
  representa la incógnita.}
\end{frame}

\begin{frame}{Ejercicio 1 — primer grado}
  \scriptsize
  \begin{alertblock}{En parejas, 12 minutos. Verifique cada respuesta.}
    \tiny
    \textbf{1.} $2x+3 = 7x-12$ \qquad
    \textbf{2.} $4(x+1)-2(x-3) = 18$ \qquad
    \textbf{3.} $\tfrac{x}{3}+\tfrac{x}{4} = 7$ \qquad
    \textbf{4.} Tres enteros consecutivos suman $129$: ¿cuáles son? \qquad
    \textbf{5.} Un rectángulo mide $7$ m más de largo que de ancho y su
    perímetro es $54$ m: halle las dimensiones.
  \end{alertblock}

  \vspace{0.05cm}
  \begin{block}{Solución}
    \tiny
    \textbf{1.} $-5x=-15 \Rightarrow x=\destacar{3}$ \qquad
    \textbf{2.} $2x+10=18 \Rightarrow x=\destacar{4}$ \qquad
    \textbf{3.} m.c.d. $12$: $4x+3x=84 \Rightarrow x=\destacar{12}$

    \vspace{0.08cm}
    \textbf{4.} $x+(x+1)+(x+2)=129 \Rightarrow 3x+3=129
    \Rightarrow x=42$: \destacar{42, 43 y 44}

    \vspace{0.08cm}
    \textbf{5.} $2(x)+2(x+7)=54 \Rightarrow 4x+14=54
    \Rightarrow x=10$: \destacar{10 m de ancho y 17 m de largo}
  \end{block}
\end{frame}

% ===========================================================================
\seccion{3.6 Ecuaciones cuadráticas}
% ===========================================================================

\begin{frame}{Tres formas de resolver una cuadrática}
  \scriptsize
  \begin{block}{Forma general: $ax^{2}+bx+c=0$, con $a\neq 0$}
    \tiny
    \textbf{Antes de nada}, iguale a cero y ordene. Si no está igualada a
    cero, ningún método funciona.
  \end{block}

  \vspace{0.04cm}
  \begin{columns}[T,onlytextwidth]
    \begin{column}{0.32\textwidth}
      \begin{exampleblock}{1. Factorizando}
        \tiny
        Si $ab=0$, entonces $a=0$ o $b=0$.

        \vspace{0.04cm}
        $x^{2}-7x+12=0$

        $(x-3)(x-4)=0$

        $x=\destacar{3}$ \ o \ $x=\destacar{4}$
      \end{exampleblock}
    \end{column}
    \begin{column}{0.32\textwidth}
      \begin{exampleblock}{2. Despejando}
        \tiny
        Cuando falta el término en $x$.

        \vspace{0.04cm}
        $2x^{2}-8=0$

        $x^{2}=4$

        $x=\destacar{\pm 2}$ \ {\color{uniGris}(aquí sí va el $\pm$)}
      \end{exampleblock}
    \end{column}
    \begin{column}{0.32\textwidth}
      \begin{exampleblock}{3. Fórmula general}
        \tiny
        Sirve siempre:

        \vspace{0.04cm}
        $x = \dfrac{-b\pm\sqrt{b^{2}-4ac}}{2a}$

        \vspace{0.04cm}
        {\color{uniGris}Se usa cuando no se ve la factorización.}
      \end{exampleblock}
    \end{column}
  \end{columns}

  \vspace{0.02cm}
  \begin{alertblock}{Si falta el término independiente, saque factor común}
    \tiny
    $x^{2}+6x = 0 \Rightarrow x(x+6)=0 \Rightarrow x=\destacar{0}$ \ o \
    $x=\destacar{-6}$. \quad
    \destacar{Nunca divida entre $x$}: perdería la solución $x=0$. Es el error
    silencioso más común del tema, porque el resultado parece razonable.
  \end{alertblock}
\end{frame}

\begin{frame}{El discriminante: cuántas soluciones hay}
  \scriptsize
  \begin{block}{$\Delta = b^{2}-4ac$}
    \tiny
    Es lo que va dentro de la raíz. Como en la sesión 8, \destacar{una raíz
    par de un negativo no existe en $\mathbb{R}$}, y de ahí sale todo.
  \end{block}

  \vspace{0.04cm}
  \begin{center}\tiny
  \begin{tabular}{@{}p{2.0cm}p{3.0cm}p{6.4cm}@{}}
    \toprule
    \textbf{Si\dots} & \textbf{Soluciones reales} & \textbf{Ejemplo} \\
    \midrule
    $\Delta > 0$ & \azul{Dos} distintas &
      $x^{2}-5x-24=0$: $\Delta=121$, \ $x=8$ o $x=-3$ \\
    $\Delta = 0$ & \azul{Una} (doble) &
      $x^{2}-10x+25=0$: $\Delta=0$, \ $x=5$ \\
    $\Delta < 0$ & \destacar{Ninguna} real &
      $x^{2}+2x+5=0$: $\Delta=-16$ \\
    \bottomrule
  \end{tabular}
  \end{center}

  \vspace{0.04cm}
  \begin{exampleblock}{La fórmula aplicada: $3x^{2}-5x-2=0$}
    \tiny
    $a=3$, $b=-5$, $c=-2$; \ $\Delta = 25+24 = 49$; \
    $x = \tfrac{5\pm 7}{6}$, luego
    $x = \destacar{2}$ \ o \ $x = \destacar{-\tfrac{1}{3}}$.
    \ \azul{Verificación con $x=2$:} $12-10-2=0$ \checkmark
  \end{exampleblock}
\end{frame}

\begin{frame}{Ejercicio 2 — cuadráticas}
  \scriptsize
  \begin{alertblock}{Individual, 12 minutos. Diga qué método usa.}
    \tiny
    \textbf{1.} $x^{2}-2x-15=0$ \qquad
    \textbf{2.} $4x^{2}-12x+9=0$ \qquad
    \textbf{3.} $2x^{2}+7x-4=0$ \qquad
    \textbf{4.} $x^{2}+6x=0$ \qquad
    \textbf{5.} $x^{2}+2x+5=0$
  \end{alertblock}

  \vspace{0.05cm}
  \begin{block}{Solución}
    \tiny
    \textbf{1.} Factorizando: $(x-5)(x+3)=0$;
    $x=\destacar{5}$ o $x=\destacar{-3}$ \qquad
    \textbf{2.} $\Delta=0$: $(2x-3)^{2}=0$;
    $x=\destacar{\tfrac{3}{2}}$ (doble)

    \vspace{0.08cm}
    \textbf{3.} Fórmula: $\Delta=49+32=81$;
    $x=\tfrac{-7\pm 9}{4}$: $x=\destacar{\tfrac{1}{2}}$ o
    $x=\destacar{-4}$ \qquad
    \textbf{4.} Factor común: $x(x+6)=0$;
    $x=\destacar{0}$ o $x=\destacar{-6}$

    \vspace{0.08cm}
    \textbf{5.} $\Delta=4-20=-16 < 0$: \destacar{no tiene solución real}
  \end{block}
\end{frame}

\begin{frame}{Problemas con cuadráticas}
  \scriptsize
  \begin{exampleblock}{Ejemplo desarrollado}
    \tiny
    Un terreno rectangular tiene $5$ m más de largo que de ancho y su área es
    $84$ m$^{2}$. Sea $x$ el ancho; entonces $x(x+5) = 84$:
    \[ x^{2}+5x-84 = 0
       \;\Longrightarrow\; (x+12)(x-7)=0
       \;\Longrightarrow\; x=-12 \ \text{o} \ x=7 \]
    \destacar{$x=-12$ se descarta}: un ancho no puede ser negativo. El terreno
    mide $\destacar{7}$ m $\times$ $\destacar{12}$ m.
  \end{exampleblock}

  \vspace{0.02cm}
  \begin{alertblock}{La lección: la ecuación da dos respuestas, el problema una}
    \tiny
    Una cuadrática casi siempre devuelve \azul{dos} soluciones matemáticas,
    pero el contexto suele admitir \destacar{solo una}. Longitudes, edades,
    cantidades de personas y precios \azul{no pueden ser negativos}.
    \destacar{Descartar con argumento es parte de la respuesta}, no un
    detalle.
  \end{alertblock}

  \vspace{0.0cm}
  {\tiny\color{uniGris} Es la misma disciplina de la sesión 11 con los
  logaritmos y de la 14 con las fracciones: \azul{resolver, y después
  verificar cuáles soluciones son admisibles}.}
\end{frame}

% ===========================================================================
\seccion{Repaso integrador de la Unidad 3}
% ===========================================================================

\begin{frame}{Mapa de la Unidad 3}
  \scriptsize
  \begin{center}
  \begin{tabular}{@{}clp{7.8cm}@{}}
    \toprule
    \textbf{Tema} & \textbf{Sesión} & \textbf{Lo que se evalúa} \\
    \midrule
    3.1 & 13 &
      Términos semejantes; sumar, restar y multiplicar polinomios \\
    \addlinespace
    3.2 & 13 &
      Los cinco productos notables y el cálculo mental \\
    \addlinespace
    3.3 & 14 &
      Factor común, diferencia de cuadrados, trinomios, agrupación, cubos;
      factorización \azul{completa} \\
    \addlinespace
    3.4 & 14 &
      Simplificar y operar fracciones algebraicas; restricciones \\
    \addlinespace
    3.5--3.6 & 15 &
      Ecuaciones de primer grado y cuadráticas; discriminante; problemas en
      contexto \\
    \bottomrule
  \end{tabular}
  \end{center}

  \vspace{0.08cm}
  {\tiny\color{uniGris} Los cinco temas están \destacar{encadenados}: sin
  productos notables no hay factorización, sin factorización no hay fracciones
  algebraicas ni cuadráticas resueltas por factoreo. Si algo del 3.3 quedó
  flojo, es lo primero que hay que repasar esta semana.}
\end{frame}

\begin{frame}{Simulacro del Examen 3 — parte A}
  \begin{alertblock}{Individual, 20 minutos, sin apuntes}
    \scriptsize
    \begin{enumerate}\setlength{\itemsep}{0.0cm}
      \item Reduzca: \ $(6x^{2}-4x+1)-(2x^{2}+3x-5)$
      \item Desarrolle: \ $(4x-3)^{2}$ \ y \ $(x+5)(x-5)$
      \item Factorice completamente: \ $2x^{3}-18x$
      \item Factorice completamente: \ $3x^{2}-12x+12$
      \item Simplifique y declare restricciones: \
            $\dfrac{x^{2}-25}{x^{2}+10x+25}$
      \item Efectúe: \ $\dfrac{2}{x-3}+\dfrac{5}{x+2}$
    \end{enumerate}
  \end{alertblock}
\end{frame}

\begin{frame}{Simulacro — solución de la parte A}
  \scriptsize
  \begin{columns}[T,onlytextwidth]
    \begin{column}{0.49\textwidth}
      \begin{block}{Puntos 1 a 3}
        \tiny
        \textbf{1.} $6x^{2}-4x+1-2x^{2}-3x+5
        = \destacar{4x^{2}-7x+6}$

        \vspace{0.1cm}
        \textbf{2.} $\destacar{16x^{2}-24x+9}$; \
        $\destacar{x^{2}-25}$

        \vspace{0.1cm}
        \textbf{3.} $2x(x^{2}-9)
        = \destacar{2x(x-3)(x+3)}$
      \end{block}
    \end{column}
    \begin{column}{0.49\textwidth}
      \begin{block}{Puntos 4 a 6}
        \tiny
        \textbf{4.} $3(x^{2}-4x+4)
        = \destacar{3(x-2)^{2}}$

        \vspace{0.1cm}
        \textbf{5.} $\dfrac{(x-5)(x+5)}{(x+5)^{2}}
        = \destacar{\dfrac{x-5}{x+5}}$, \ $x\neq -5$

        \vspace{0.1cm}
        \textbf{6.} $\dfrac{2(x+2)+5(x-3)}{(x-3)(x+2)}
        = \destacar{\dfrac{7x-11}{x^{2}-x-6}}$
      \end{block}
    \end{column}
  \end{columns}

  \vspace{0.08cm}
  {\tiny\color{uniGris} En el \destacar{3} y el \destacar{4}, quedarse en
  $2x(x^{2}-9)$ o $3(x^{2}-4x+4)$ es factorización \azul{incompleta}: vale
  media respuesta. En el \destacar{5}, olvidar la restricción $x\neq -5$ es el
  descuento más común de todo el examen.}
\end{frame}

\begin{frame}{Simulacro del Examen 3 — parte B}
  \begin{alertblock}{Individual, 20 minutos}
    \scriptsize
    \begin{enumerate}\setlength{\itemsep}{0.0cm}
      \setcounter{enumi}{6}
      \item Resuelva: \ $3x-7 = 14$
      \item Resuelva: \ $\dfrac{x-1}{2}+\dfrac{x+3}{5} = 4$
      \item Resuelva factorizando: \ $x^{2}-7x+12 = 0$
      \item Resuelva con la fórmula: \ $2x^{2}+7x-4 = 0$
      \item Sin resolver, diga cuántas soluciones reales tiene
            $x^{2}+2x+5=0$ y por qué.
      \item Un terreno rectangular es $5$ m más largo que ancho y su área es
            $84$ m$^{2}$. Halle las dimensiones.
    \end{enumerate}
  \end{alertblock}
\end{frame}

\begin{frame}{Simulacro — solución de la parte B}
  \scriptsize
  \begin{columns}[T,onlytextwidth]
    \begin{column}{0.49\textwidth}
      \begin{block}{Puntos 7 a 9}
        \tiny
        \textbf{7.} $3x=21 \Rightarrow x=\destacar{7}$

        \vspace{0.1cm}
        \textbf{8.} Por $10$: $5(x-1)+2(x+3)=40$;
        $7x+1=40 \Rightarrow x=\destacar{\tfrac{39}{7}}$

        \vspace{0.1cm}
        \textbf{9.} $(x-3)(x-4)=0$;
        $x=\destacar{3}$ o $x=\destacar{4}$
      \end{block}
    \end{column}
    \begin{column}{0.49\textwidth}
      \begin{block}{Puntos 10 a 12}
        \tiny
        \textbf{10.} $\Delta=49+32=81$; $x=\tfrac{-7\pm 9}{4}$:
        $\destacar{\tfrac{1}{2}}$ y $\destacar{-4}$

        \vspace{0.1cm}
        \textbf{11.} $\Delta = 4-20 = -16 < 0$:
        \destacar{ninguna} real

        \vspace{0.1cm}
        \textbf{12.} $x(x+5)=84 \Rightarrow (x+12)(x-7)=0$;
        se descarta $-12$: \destacar{7 m $\times$ 12 m}
      \end{block}
    \end{column}
  \end{columns}

  \vspace{0.08cm}
  {\tiny\color{uniGris} \azul{Verificaciones que caben en el margen:} en el
  \destacar{8}, $\tfrac{39}{7}\approx 5{,}571$ y
  $\tfrac{4{,}571}{2}+\tfrac{8{,}571}{5} = 2{,}286+1{,}714 = 4$ \checkmark; en
  el \destacar{10}, $2(\tfrac14)+7(\tfrac12)-4 = 0{,}5+3{,}5-4 = 0$
  \checkmark}
\end{frame}

\begin{frame}{Errores frecuentes de toda la Unidad 3}
  \scriptsize
  \begin{center}
  \begin{tabular}{@{}p{5.2cm}p{6.8cm}@{}}
    \toprule
    \textbf{Error} & \textbf{Cómo evitarlo} \\
    \midrule
    $(a+b)^{2}=a^{2}+b^{2}$ &
      Falta $2ab$. Compruebe con $a=b=1$ \\
    \addlinespace
    Cancelar sumandos: $\tfrac{x+3}{x}=3$ &
      Solo se cancelan \azul{factores}; factorice primero \\
    \addlinespace
    Factorización incompleta &
      Pregúntese si cada factor todavía factoriza \\
    \addlinespace
    No declarar restricciones &
      Anótelas \azul{antes} de simplificar \\
    \addlinespace
    Dividir entre $x$ en $x^{2}+6x=0$ &
      Se pierde $x=0$: saque factor común \\
    \addlinespace
    Resolver sin igualar a cero &
      La forma $ax^{2}+bx+c=0$ es obligatoria \\
    \addlinespace
    Dar las dos raíces en un problema de área &
      Descarte la negativa \azul{con argumento} \\
    \bottomrule
  \end{tabular}
  \end{center}
\end{frame}

\begin{frame}{Cierre del curso}
  \scriptsize
  \begin{block}{Lo que se lleva de Matemáticas Operativas}
    \tiny
    \textbf{Unidad 1:} los campos numéricos y cómo operar con ellos sin
    calculadora, más razones, proporciones y porcentajes. \quad
    \textbf{Unidad 2:} potencias, radicales, racionalización y logaritmos.
    \quad
    \textbf{Unidad 3:} el lenguaje algebraico completo y la resolución de
    ecuaciones.

    \vspace{0.06cm}
    Todo esto es el \destacar{prerrequisito operativo} de los cursos que
    siguen —Cálculo, Estadística, Finanzas—: allí no se vuelve a explicar cómo
    factorizar ni cómo despejar, se da por sabido.
  \end{block}

  \vspace{0.1cm}
  \begin{alertblock}{Para preparar el Examen 3 (30\%) — semana 16}
    \tiny
    \textbf{1.} Rehaga el simulacro de hoy \azul{cronometrado} y sin mirar las
    soluciones. \quad
    \textbf{2.} Repase los seis casos de factorización de la sesión 14: son la
    base de casi todos los puntos. \quad
    \textbf{3.} Memorice los cinco productos notables y la fórmula general.
    \quad
    \textbf{4.} En cada ejercicio, \destacar{verifique} reemplazando.
  \end{alertblock}

  \vspace{0.08cm}
  \begin{center}
    {\footnotesize\color{uniAzul}\textbf{Gracias por el semestre.}
    \ \ Consultas y asesorías: hasta el día anterior al examen.}
  \end{center}
\end{frame}

\begin{frame}{Ejercicios propuestos}
  \vspace{-0.1cm}
  \begin{beamercolorbox}[wd=\textwidth,sep=0.10cm]{block body}
    \centering\scriptsize
    \textbf{1--7:} primer grado. \quad \textbf{8--15:} cuadráticas. \quad
    \textbf{16--20:} problemas.
  \end{beamercolorbox}
  \listatiny
  \vspace{0.15cm}
  \begin{columns}[T,onlytextwidth]
    \begin{column}{0.47\textwidth}
      \begin{enumerate}\setlength{\itemsep}{0.16cm}
        \item $7x-4 = 3x+16$
        \item $5(x-3) = 2(x+6)$
        \item $\tfrac{x}{2}-\tfrac{x}{5} = 3$
        \item $\tfrac{2x+1}{3} = \tfrac{x-2}{2}$
        \item $3(x-1)-4(x+2) = -15$
        \item $0{,}25x + 0{,}1(x-40) = 45$
        \item $\tfrac{x}{4}+\tfrac{x+4}{6} = 4$
        \item $x^{2}-9x+18 = 0$
        \item $x^{2}+4x-21 = 0$
        \item $x^{2}-49 = 0$
      \end{enumerate}
    \end{column}
    \begin{column}{0.50\textwidth}
      \begin{enumerate}\setlength{\itemsep}{0.16cm}
        \setcounter{enumi}{10}
        \item $x^{2}-8x = 0$
        \item $x^{2}+6x+9 = 0$
        \item $2x^{2}-5x-3 = 0$
        \item $3x^{2}+2x+5 = 0$
        \item $x^{2}-4x+1 = 0$ \ (deje la raíz indicada)
        \item Tres enteros consecutivos suman $174$
        \item Un número más su triple es $92$
        \item Un rectángulo es $3$ m más largo que ancho y su área es
              $70$ m$^{2}$
        \item Con IVA del $19\%$ un artículo cuesta $\$178\,500$: halle la
              base
        \item Dos números suman $20$ y su producto es $91$
      \end{enumerate}
    \end{column}
  \end{columns}
\end{frame}

\begin{frame}{Respuestas}{Autoevaluación de los ejercicios propuestos}
  \listatiny
  \vspace{0.15cm}
  \begin{columns}[T,onlytextwidth]
    \begin{column}{0.47\textwidth}
      \begin{enumerate}\setlength{\itemsep}{0.16cm}
        \item $x=5$
        \item $x=9$
        \item $x=10$
        \item $x=-8$
        \item $x=4$
        \item $x=140$
        \item $x=8$
        \item $x=3$ o $x=6$
        \item $x=3$ o $x=-7$
        \item $x=\pm 7$
      \end{enumerate}
    \end{column}
    \begin{column}{0.50\textwidth}
      \begin{enumerate}\setlength{\itemsep}{0.16cm}
        \setcounter{enumi}{10}
        \item $x=0$ o $x=8$
        \item $x=-3$ (doble)
        \item $x=3$ o $x=-\tfrac{1}{2}$
        \item Ninguna real ($\Delta=-56$)
        \item $x=2\pm\sqrt{3}$
        \item $57,\ 58,\ 59$
        \item $23$
        \item $7$ m $\times$ $10$ m
        \item $\$150\,000$
        \item $7$ y $13$
      \end{enumerate}
    \end{column}
  \end{columns}
\end{frame}

\begin{frame}{Referencias}
  \footnotesize
  \begin{itemize}
    \item[{[1]}] Ramos del Olmo, Á. M. (2016). \textit{Matemáticas básicas
      para el acceso a la universidad}. Ediciones Pirámide.
      \texttt{https://elibro.net/es/ereader/uniminuto/49134}
    \item[{[2]}] González Pulido, J. W., Chacón Benavides, J. A. \&
      Fonseca Correa, L. Á. (2023). \textit{Matemática básica}
      (1.\textsuperscript{a} ed.). Universidad Mariana.
      \texttt{https://elibro.net/es/ereader/uniminuto/280647}
    \item[{[3]}] Torrijos, M. (2018). \textit{Apuntes de matemática básica}.
      Universidad Católica de Colombia.
      \texttt{https://elibro.net/es/ereader/uniminuto/197048}
    \item[{[4]}] Ramírez, A. \& Rojas, L. (2016). \textit{Matemáticas básicas
      con aplicaciones a la ingeniería}. Ecoe Ediciones.
      \texttt{https://elibro.net/es/ereader/uniminuto/114355}
  \end{itemize}

  \vspace{0.3cm}
  \begin{center}
    {\Large\color{uniAzul}\textbf{¿Preguntas?}}
  \end{center}
\end{frame}

\end{document}
